\documentclass[10pt]{article}
\usepackage{graphicx}
\setlength{\textwidth}{6.85in}
\setlength{\oddsidemargin}{-0.18in}
\setlength{\evensidemargin}{-0.18in}
\setlength{\textheight}{9.45in}
\setlength{\topmargin}{-0.65in}
\setlength{\parindent}{0pt}
\setlength{\parskip}{0.4em}
\setlength{\emergencystretch}{3em}
\title{\textbf{Wind Direction Convention Error in the Downwind/Upwind Construction}}
\author{Reproducible diagnostic for the bureaucrats and farms project}
\date{August 2026}

\begin{document}
\maketitle

\section*{Executive finding}
The legacy pipeline combines two angular coordinate systems without converting
between them. The wind variable is constructed with
$\mathrm{atan2}(v,u)$, which is a mathematical angle measured
counterclockwise from \textbf{East}. The geographic functions
\texttt{geosphere::bearing()} and \texttt{geosphere::destPoint()} interpret
their angles as compass bearings measured clockwise from \textbf{North}. Both values are
expressed in degrees, but they are not directly comparable.

This is not a cosmetic labeling issue. In the second scenario below, the true
wind points northwest and the AC representative point lies south-southeast of
the grid, so the representative point is unambiguously upwind. The legacy code
instead interprets the same numeric value as a southeast-pointing wind and
classifies the representative point as downwind.

\section*{A real project example}
The script searches the project shapefiles rather than inventing coordinates.
It selects grid \textbf{116142} in the
\textbf{SUNAM} assembly constituency
(AC ID \textbf{449}), SANGRUR, Punjab. The grid centroid
is 9.96 km from the AC representative point. The compass
bearing from the grid centroid to that point is
\textbf{164.986$^\circ$}, effectively the requested
165$^\circ$ example.

For a mathematical angle $\alpha$ measured from East, the corresponding
clockwise-from-North direction toward which the wind travels is
\[
  B_{\mathrm{toward}} = (90^\circ - \alpha) \bmod 360^\circ.
\]
The legacy code instead uses the numeric value $\alpha$ directly as if it were
a compass bearing.

\begin{table}[ht]
\centering
\small
\resizebox{\textwidth}{!}{%
\begin{tabular}{rrrrrrll}
\hline
$\alpha$ from East & Correct compass & Legacy compass & AC bearing &
$\Delta$ correct & $\Delta$ legacy & Correct & Legacy \\
\hline
55 & 35 & 55 & 165.0 & 130 & 110 & Upwind & Upwind \\
135 & 315 & 135 & 165.0 & 150 & 30 & Upwind & \textbf{Downwind (WRONG)} \\
\hline
\end{tabular}
}
\caption{Angles and classifications. A point is downwind only when its circular
angular separation from the wind-toward direction is less than 90$^\circ$.}
\end{table}

\clearpage
\begin{figure}[p]
\centering
\includegraphics[width=\textwidth]{wind\_direction\_convention\_error\_figure.pdf}
\caption{Four views of the same real AC/grid pair. Blue panels use the correct
conversion from an East-referenced mathematical angle to a compass bearing.
Red panels reproduce the legacy interpretation. The shaded portion is the AC
half-plane labeled downwind; the black arrow points from the grid to the AC
representative point. With a 55$^\circ$ input, both procedures call the point
upwind, but only by chance. With a 135$^\circ$ input, the correct wind points
northwest while the legacy interpretation points southeast, producing the
incorrect downwind classification.}
\end{figure}
\clearpage

\section*{What is actually stored in \texttt{downup\_ac\_area}}
The area indicator is defined as 1 when the area stored as downwind is larger
than the area stored as upwind, and 0 otherwise. The diagnostic script reads
the actual area parquet, selects the same real grid, and recomputes both
half-plane intersections from the shapefile. Treating the raw wind number as a
North-referenced compass bearing reproduces the stored areas to within
$10^{-6}$ km$^2$. This proves that the convention error is present in the
saved data, not only in a hypothetical reconstruction.

\begin{table}[ht]
\centering
\scriptsize
\resizebox{\textwidth}{!}{%
\begin{tabular}{rrrrrrrrr}
\hline
Period & Wind from East & Stored down & Stored up & Stored flag &
Correct compass & Correct down & Correct up & Correct flag \\
\hline
2017-07 & 119.4 & 459.7 & 104.7 & \textbf{1 (WRONG)} & 330.6 & 24.0 & 540.4 & \textbf{0} \\
2017-03 & -75.0 & 159.7 & 404.7 & \textbf{0 (WRONG)} & 165.0 & 557.2 & 7.2 & \textbf{1} \\
\hline
\end{tabular}
}
\caption{Actual saved areas and classifications versus exact corrected
intersections. Areas are in km$^2$; \texttt{downup\_ac\_area} equals 1
when downwind area is larger than upwind area.}
\end{table}

For this selected grid, 136 of
148 nonmissing monthly observations change classification.
The complete stored-by-corrected count is:
\begin{center}
\begin{tabular}{lrr}
\hline
 & Corrected 0 & Corrected 1 \\
\hline
Stored 0 & 12 &
112 \\
Stored 1 & 24 &
0 \\
\hline
\end{tabular}
\end{center}
These counts describe this one diagnostic grid only; they are not an estimate
of the error rate in the full sample.

\section*{Why the negative angle is also misclassified}
In March 2017, the stored raw direction is
-74.971$^\circ$ from East. A negative
Cartesian angle is valid: it points clockwise below the East axis. The correct
compass-toward direction is
$(90-(-74.971))\bmod 360=
164.971^\circ$. The legacy code
instead normalizes the unconverted number to
285.029$^\circ$ and reads it from
North. It therefore stores \texttt{downup\_ac\_area} =
0, whereas the corrected majority-area
classification is 1.

\clearpage
\begin{figure}[p]
\centering
\includegraphics[width=\textwidth]{wind\_direction\_convention\_error\_actual\_area\_figure.pdf}
\caption{The values actually stored in the parquet (left column) versus
the corrected calculations (right column). The first row is a positive wind
angle; the second row is a negative wind angle. In every panel the black arrow
to the AC representative point is a compass bearing measured clockwise from
North. The wind input printed at the top is a mathematical angle measured from
East. Red shading is the area stored as downwind; blue shading is the corrected
downwind area. Both examples reverse \texttt{downup\_ac\_area}.}
\end{figure}
\clearpage

\section*{Corrected two-square area method}
The corrected method first converts the wind angle to the compass-toward
direction $B$. At the focal grid centroid $p$, define the unit wind vector
$w=(\sin B,\cos B)$ and its perpendicular tangent
$t=(\cos B,-\sin B)$. A large downwind square has corners

\[
p-Rt,\quad p+Rt,\quad p+Rt+2Rw,\quad p-Rt+2Rw,
\]

and the upwind square is constructed in the opposite direction, using $-w$.
The radius $R$ is chosen larger than the maximum distance from the grid
centroid to the AC bounds. Therefore, the two squares cover the two complete
halves of the AC and share only the perpendicular dividing line.

The exact areas are then
\[
A_{down}=\mathrm{area}(AC\cap S_{down}),\qquad
A_{up}=\mathrm{area}(AC\cap S_{up}),\qquad
\texttt{downup\_ac\_area}=\mathbf{1}[A_{down}>A_{up}].
\]
This construction is equivalent to clipping the AC by complementary
half-planes, but showing the two finite squares makes the implementation and
area accounting explicit.

\begin{figure}[ht]
\centering
\includegraphics[width=\textwidth]{wind\_direction\_convention\_error\_square\_method\_figure.pdf}
\caption{The corrected two-square construction for the July 2017 example.
Panel (a) shows the complementary squares and their shared line, which is
perpendicular to the corrected wind vector. Panels (b) and (c) show the exact
AC intersections whose areas determine the binary classification.}
\end{figure}
\clearpage

\section*{Where the error enters}
The wind-cleaning code calculates
\begin{verbatim}
wind_direction = atan2(v10, u10) * 180 / pi
\end{verbatim}
This returns a Cartesian angle: 0$^\circ$ is East, 90$^\circ$ is North,
180$^\circ$ is West, and -90$^\circ$ is South. The downstream R code then uses
\begin{verbatim}
angle1 = wind_direction + 90
angle2 = wind_direction - 90
destPoint(grid_centroid, angle1, ...)
bearing_relative = bearing(grid_centroid, ac_point)
\end{verbatim}
Here, 0$^\circ$ is North and 90$^\circ$ is East. The subsequent numeric
comparison therefore mixes an East-referenced angle with a North-referenced
bearing. In addition, ordinary inequalities do not correctly handle intervals
that cross 0$^\circ$/360$^\circ$.

\section*{Required correction}
The wind should first be represented as a compass direction toward which the
air moves:
\begin{verbatim}
wind_direction_to = (atan2(u10, v10) * 180 / pi + 360) %% 360
\end{verbatim}
The downwind test should use circular angular distance:
\begin{verbatim}
angle_difference = ((bearing_relative - wind_direction_to + 180) %% 360) - 180
downwind = as.integer(abs(angle_difference) < 90)
\end{verbatim}
Monthly and rolling directions should also use circular means, preferably by
aggregating vector components, rather than arithmetic means of degree values.
For example, the arithmetic mean of 359$^\circ$ and 1$^\circ$ is 180$^\circ$,
while their circular mean is 0$^\circ$.

\section*{Implications for existing outputs}
\begin{itemize}
  \item The population and area routines can agree with each other while both
        label the wrong geographical half-plane.
  \item Existing \texttt{downup\_ac\_pop} and \texttt{downup\_ac\_area}
        values should not be treated as directionally validated.
  \item Corrected wind directions require regenerating the wind panel,
        population measures, area measures, master dataset, and stacked data.
  \item Cardinal and diagonal unit tests should be retained so that East,
        North, West, South, 55$^\circ$, and 135$^\circ$ cannot silently change
        convention again.
\end{itemize}

\section*{Source documentation}
ECMWF defines positive $u$ as west-to-east flow and positive $v$ as
south-to-north flow. See
\begin{flushleft}\small\ttfamily
\detokenize{https://confluence.ecmwf.int/spaces/FCST/pages/111155337/}
\end{flushleft}
ECMWF's discussion of meteorological wind angles and the special care required
with \texttt{atan2} is available at
\begin{flushleft}\small\ttfamily
\detokenize{https://confluence.ecmwf.int/pages/viewpage.action?pageId=226500971}
\end{flushleft}

\end{document}
